% Graphs and plots from the source deck were intentionally omitted or replaced
% with placeholders, per user request.
\documentclass[aspectratio=43,9pt]{beamer}

\usepackage[T1]{fontenc}
\usepackage[utf8]{inputenc}
\usepackage{lmodern}
\usepackage{amsmath,amssymb,booktabs,tabularx,array,multirow}
\usepackage{listings}
\usepackage{xcolor}
\usepackage{graphicx}
\usepackage{ragged2e}

\usetheme{default}
\setbeamertemplate{navigation symbols}{}
\setbeamertemplate{footline}[frame number]
\setbeamersize{text margin left=0.35cm,text margin right=0.35cm}
\setbeamercolor{frametitle}{fg=black,bg=white}
\setbeamerfont{frametitle}{series=\bfseries,size*= {18pt}{20pt}}
\setbeamercolor{title}{fg=black}
\setbeamerfont{title}{series=\bfseries,size*= {25pt}{27pt}}
\setbeamercolor{block title}{bg=black,fg=white}
\setbeamercolor{block body}{bg=black!5,fg=black}
\setbeamercolor{alertblock title}{bg=black,fg=white}
\setbeamercolor{alertblock body}{bg=black!5,fg=black}
\setbeamercolor{exampleblock title}{bg=black,fg=white}
\setbeamercolor{exampleblock body}{bg=black!5,fg=black}
\setbeamercolor{structure}{fg=black}
\setbeamercolor{itemize item}{fg=black}
\setbeamercolor{itemize subitem}{fg=black}
\setbeamercolor{itemize subsubitem}{fg=black}
\setbeamercolor{alerted text}{fg=black}
\setbeamercolor{example text}{fg=black}
\setbeamertemplate{itemize item}{\color{black}\small\textbullet}
\setbeamertemplate{itemize subitem}{\color{black}\footnotesize\textbullet}
\setbeamertemplate{itemize subsubitem}{\color{black}\tiny\textbullet}
\setbeamerfont{block title}{series=\bfseries}
\setlength{\tabcolsep}{3pt}
\renewcommand{\arraystretch}{1.08}
\hypersetup{colorlinks=true,linkcolor=black,urlcolor=black,breaklinks=true}

\lstset{
  basicstyle=\ttfamily\tiny,
  breaklines=true,
  columns=fullflexible,
  keepspaces=true,
  showstringspaces=false,
  frame=single,
  framerule=0.2pt,
  xleftmargin=0pt,
  xrightmargin=0pt,
  aboveskip=2pt,
  belowskip=2pt
}

\newcommand{\graphplaceholder}[3][\linewidth]{%
  \fbox{\begin{minipage}[c][#2][c]{#1}\centering\scriptsize #3\end{minipage}}%
}

\newcommand{\srclink}[2]{\href{#2}{\textcolor{black!68}{#1}}}
\newcommand{\sep}{;\hspace{0.35em}}
\newcommand{\refline}[2][References]{%
  \vfill
  {\raggedright\fontsize{5.1}{5.5}\selectfont\color{black!68}\textbf{#1:} #2\par}%
}
\newcommand{\srcChen}{\srclink{M. Chen et al., ``Evaluating large language models trained on code,'' arXiv (2021)}{https://arxiv.org/abs/2107.03374}}
\newcommand{\srcVaswani}{\srclink{A. Vaswani et al., ``Attention Is All You Need,'' NeurIPS (2017)}{https://arxiv.org/abs/1706.03762}}
\newcommand{\srcRadford}{\srclink{A. Radford et al., ``Improving Language Understanding by Generative Pre-Training'' (2018)}{https://cdn.openai.com/research-covers/language-unsupervised/language_understanding_paper.pdf}}
\newcommand{\srcDevlin}{\srclink{J. Devlin et al., ``BERT: Pre-training of Deep Bidirectional Transformers for Language Understanding,'' NAACL-HLT (2019)}{https://aclanthology.org/N19-1423/}}
\newcommand{\srcChild}{\srclink{R. Child et al., ``Generating Long Sequences with Sparse Transformers,'' arXiv (2019)}{https://arxiv.org/abs/1904.10509}}
\newcommand{\srcHusain}{\srclink{H. Husain et al., ``CodeSearchNet Challenge,'' arXiv (2019)}{https://arxiv.org/abs/1909.09436}}
\newcommand{\srcGao}{\srclink{L. Gao et al., ``The Pile,'' arXiv (2020)}{https://arxiv.org/abs/2101.00027}}
\newcommand{\srcFeng}{\srclink{Z. Feng et al., ``CodeBERT,'' Findings of EMNLP (2020)}{https://arxiv.org/abs/2002.08155}}
\newcommand{\srcBrown}{\srclink{T. Brown et al., ``Language Models are Few-Shot Learners,'' NeurIPS (2020)}{https://arxiv.org/abs/2005.14165}}
\newcommand{\srcWang}{\srclink{B. Wang and A. Komatsuzaki, ``GPT-J-6B'' (2021)}{https://arankomatsuzaki.wordpress.com/2021/06/04/gpt-j/}}
\newcommand{\srcPapineni}{\srclink{K. Papineni et al., ``BLEU,'' ACL (2002)}{https://aclanthology.org/P02-1040/}}
\newcommand{\srcRen}{\srclink{S. Ren et al., ``CodeBLEU,'' arXiv (2020)}{https://arxiv.org/abs/2009.10297}}
\newcommand{\srcRoziere}{\srclink{B. Roziere et al., ``Unsupervised Translation of Programming Languages,'' NeurIPS (2020)}{https://arxiv.org/abs/2006.03511}}
\newcommand{\srcKulal}{\srclink{S. Kulal et al., ``SPoC: Search-based Pseudocode to Code,'' NeurIPS (2019)}{https://arxiv.org/abs/1906.04908}}
\newcommand{\srcHumanEval}{\srclink{OpenAI HumanEval benchmark repository}{https://github.com/openai/human-eval}}
\newcommand{\srcGVisor}{\srclink{gVisor documentation}{https://gvisor.dev/}}
\newcommand{\srcKingmaBa}{\srclink{D. Kingma and J. Ba, ``Adam,'' ICLR (2015)}{https://arxiv.org/abs/1412.6980}}
\newcommand{\srcHoltzman}{\srclink{A. Holtzman et al., ``The Curious Case of Neural Text Degeneration,'' ICLR (2020)}{https://arxiv.org/abs/1904.09751}}
\newcommand{\srcKaplan}{\srclink{J. Kaplan et al., ``Scaling Laws for Neural Language Models,'' arXiv (2020)}{https://arxiv.org/abs/2001.08361}}
\newcommand{\srcHendrycks}{\srclink{D. Hendrycks et al., ``Measuring Coding Challenge Competence with APPS,'' NeurIPS (2021)}{https://arxiv.org/abs/2105.09938}}
\newcommand{\srcBlack}{\srclink{S. Black et al., ``GPT-Neo'' project release (2021)}{https://zenodo.org/records/5297715}}
\newcommand{\srcPrime}{\srclink{Wikipedia, ``Primality test''}{https://en.wikipedia.org/wiki/Primality_test}}
\newcommand{\srcTravis}{\srclink{Travis CI documentation}{https://www.travis-ci.com/}}
\newcommand{\srcTox}{\srclink{Tox documentation}{https://tox.wiki/}}
\newcommand{\srcMcCabe}{\srclink{T. McCabe, ``A Complexity Measure,'' IEEE TSE (1976)}{https://www.computer.org/csdl/journal/ts/1976/04/01702388/13rRUxYIN5N}}
\newcommand{\srcONeill}{\srclink{M. O'Neill et al., ``Automatic Programming: The Open Issue?'' (2020)}{https://www.nber.org/papers/w25682}}
\newcommand{\srcClarkson}{\srclink{M. Clarkson and F. Schneider, ``Hyperproperties,'' CSF (2010)}{https://arxiv.org/abs/1401.4492}}
\newcommand{\srcLeveson}{\srclink{N. Leveson, ``Improving the Standard Risk Matrix'' (2019)}{https://sunnyday.mit.edu/Risk-Matrix.pdf}}
\newcommand{\srcChristiano}{\srclink{P. Christiano, ``Clarifying AI Alignment'' (2018)}{https://ai-alignment.com/clarifying-ai-alignment-cec47cd69dd6}}
\newcommand{\srcKenton}{\srclink{Z. Kenton et al., ``Alignment of Language Agents'' (2021)}{https://arxiv.org/abs/2103.14659}}
\newcommand{\srcKeskar}{\srclink{N. Keskar et al., ``CTRL,'' arXiv (2019)}{https://arxiv.org/abs/1909.05858}}
\newcommand{\srcStiennon}{\srclink{N. Stiennon et al., ``Learning to Summarize with Human Feedback,'' NeurIPS (2020)}{https://arxiv.org/abs/2009.01325}}
\newcommand{\srcBlodgett}{\srclink{S. Blodgett et al., ``Language (Technology) Is Power,'' ACL (2020)}{https://aclanthology.org/2020.acl-main.485/}}
\newcommand{\srcBender}{\srclink{E. Bender et al., ``On the Dangers of Stochastic Parrots,'' FAccT (2021)}{https://s10251.pcdn.co/pdf/2021-bender-parrots.pdf}}
\newcommand{\srcAbid}{\srclink{A. Abid et al., ``Persistent Anti-Muslim Bias in Large Language Models,'' AIES (2021)}{https://arxiv.org/abs/2101.05783}}
\newcommand{\srcCrawfordBias}{\srclink{K. Crawford, ``The Trouble with Bias'' (2017)}{https://ainowinstitute.org/news/the-trouble-with-bias}}
\newcommand{\srcAcemoglu}{\srclink{D. Acemoglu, ``The Wrong Kind of AI?'' CJRES (2021)}{https://academic.oup.com/cjres/article/14/1/25/6153095}}
\newcommand{\srcCarlini}{\srclink{N. Carlini et al., ``Extracting Training Data from Large Language Models,'' USENIX Security (2021)}{https://www.usenix.org/conference/usenixsecurity21/presentation/carlini-extracting-training-data}}
\newcommand{\srcSchwartz}{\srclink{R. Schwartz et al., ``Green AI,'' CACM (2020)}{https://arxiv.org/abs/1907.10597}}
\newcommand{\srcAmodei}{\srclink{D. Amodei and D. Hernandez, ``AI and Compute'' (2018)}{https://openai.com/index/ai-and-compute/}}
\newcommand{\srcSmith}{\srclink{B. Smith, ``Microsoft Will Be Carbon Negative by 2030'' (2020)}{https://blogs.microsoft.com/blog/2020/01/16/microsoft-will-be-carbon-negative-by-2030/}}
\newcommand{\srcCrawfordAtlas}{\srclink{K. Crawford, The Atlas of AI (2021)}{https://yalebooks.yale.edu/book/9780300264630/atlas-of-ai/}}
\newcommand{\srcZiegler}{\srclink{A. Ziegler, ``A First Look at Rote Learning in GitHub Copilot Suggestions'' (2021)}{https://github.blog/ai-and-ml/github-copilot/github-copilot-research-recitation/}}
\newcommand{\srcZaremba}{\srclink{W. Zaremba and I. Sutskever, ``Learning to Execute,'' arXiv (2014)}{https://arxiv.org/abs/1410.4615}}
\newcommand{\srcHindle}{\srclink{A. Hindle et al., ``On the Naturalness of Software,'' ICSE (2012)}{https://softwareprocess.es/homepage/papers/2012-hindle12012icse/}}
\newcommand{\srcCopilotMSR}{\srclink{P. Mozannar et al., ``Research: Quantifying GitHub Copilot's Impact on Developer Productivity and Happiness'' (2023)}{https://arxiv.org/abs/2302.06590}}
\newcommand{\srcCui}{\srclink{Z. Cui et al., ``The Effects of Generative AI on High-Skilled Work,'' Microsoft Research (2025)}{https://www.microsoft.com/en-us/research/publication/the-effects-of-generative-ai-on-high-skilled-work-evidence-from-three-field-experiments-with-software-developers/}}
\newcommand{\srcMETR}{\srclink{METR, ``Measuring the Impact of Early-2025 AI on Experienced Open-Source Developer Productivity'' (2025)}{https://arxiv.org/abs/2507.09089}}
\newcommand{\srcILO}{\srclink{ILO, ``Generative AI and Jobs: 2025 Update''}{https://www.ilo.org/publications/generative-ai-and-jobs-2025-update}}
\newcommand{\srcBLS}{\srclink{U.S. BLS, ``AI Impacts in BLS Employment Projections'' (2025)}{https://www.bls.gov/opub/ted/2025/ai-impacts-in-bls-employment-projections.htm}}
\newcommand{\srcWEF}{\srclink{World Economic Forum, ``Future of Jobs Report 2025''}{https://www.weforum.org/publications/the-future-of-jobs-report-2025/digest/}}
\newcommand{\srcSWEBench}{\srclink{SWE-bench benchmark repository}{https://github.com/swe-bench/SWE-bench}}
\newcommand{\srcSWERebench}{\srclink{SWE-rebench leaderboard and reports}{https://swe-rebench.com/}}
\newcommand{\srcMHumanEval}{\srclink{mHumanEval benchmark (NAACL 2025)}{https://aclanthology.org/2025.naacl-long.570/}}
\newcommand{\srcIEA}{\srclink{IEA, Energy and AI (2025)}{https://www.iea.org/reports/energy-and-ai/}}
\newcommand{\srcUSCOPTwo}{\srclink{U.S. Copyright Office, Copyright and AI Part 2: Copyrightability (2025)}{https://www.copyright.gov/ai/Copyright-and-Artificial-Intelligence-Part-2-Copyrightability-Report.pdf}}
\newcommand{\srcUSCOPThree}{\srclink{U.S. Copyright Office, Copyright and AI Part 3: Generative AI Training (2025)}{https://www.copyright.gov/ai/Copyright-and-Artificial-Intelligence-Part-3-Generative-AI-Training-Report-Pre-Publication-Version.pdf}}
\newcommand{\srcEUGPAI}{\srclink{European Commission, GPAI training-data transparency template (2025)}{https://digital-strategy.ec.europa.eu/en/news/commission-presents-template-general-purpose-ai-model-providers-summarise-data-used-train-their}}
\newcommand{\srcNIST}{\srclink{NIST, AI 600-1: Managing Misuse Risk for Dual-Use Foundation Models (2025)}{https://nvlpubs.nist.gov/nistpubs/ai/NIST.AI.600-1.pdf}}
\newcommand{\srcPromptInj}{\srclink{``What Is the Threat from Prompt Injection to Agentic Systems?'' (2026)}{https://arxiv.org/html/2601.17548v1}}
\newcommand{\srcOWASP}{\srclink{OWASP Top 10 for LLM Applications 2025}{https://owasp.org/www-project-top-10-for-large-language-model-applications/assets/PDF/OWASP-Top-10-for-LLMs-v2025.pdf}}

\newcommand{\outlineitem}[2]{%
  \ifnum\currentstage=#1\relax
    \item[{\textcolor{black}{\large\textbullet}}] {\color{black}\bfseries #2}
  \else
    \item[{\textcolor{black}{\large\textbullet}}] {\color{black!70}#2}
  \fi
}

\newcommand{\outlineframe}[1]{%
\begin{frame}{Outline}
\small
\def\currentstage{#1}
\begin{itemize}
  \setlength{\itemsep}{0.45em}
  \outlineitem{1}{Background and approach}
  \outlineitem{2}{Evaluation}
  \outlineitem{3}{Code fine-tuning}
  \outlineitem{4}{Experiments}
  \outlineitem{5}{Supervised fine-tuning}
  \outlineitem{6}{Docstring generation}
  \outlineitem{7}{Limitations}
  \outlineitem{8}{Broader impacts}
  \outlineitem{9}{Related work}
\end{itemize}
\end{frame}
}

\begin{document}

% 1
\begin{frame}[plain]
\centering
\vspace*{0.35em}
{\usebeamerfont{title}\bfseries Evaluating Large Language\\[0.18em]Models Trained on Code (Codex)\par}
\vspace{0.9em}
{\scriptsize
M. Chen, J. Tworek, H. Jun, Q. Yuan, H. Ponde de Oliveira Pinto, J. Kaplan, H. Edwards, Y. Burda, N. Joseph, G. Brockman, A. Ray, R. Puri, G. Krueger, M. Petrov, H. Khlaaf, G. Sastry, P. Mishkin, B. Chan, S. Gray, N. Ryder, M. Pavlov, A. Power, L. Kaiser, M. Bavarian, C. Winter, P. Tillet, F. Petroski Such, D. Cummings, M. Plappert, F. Chantzis, E. Barnes, A. Herbert-Voss, W. Hebgen Guss, A. Nichol, A. Paino, N. Tezak, J. Tang, I. Babuschkin, S. Balaji, S. Jain, W. Saunders, C. Hesse, A. N. Carr, J. Leike, J. Achiam, V. Misra, E. Morikawa, A. Radford, M. Knight, M. Brundage, M. Murati, K. Mayer, P. Welinder, B. McGrew, D. Amodei, S. McCandlish, I. Sutskever, W. Zaremba\par}
\vspace{0.35em}
{\small\href{https://arxiv.org/abs/2107.03374}{arXiv (2021)}\par}
\end{frame}

% 2
\outlineframe{1}

% 3
\begin{frame}{Background}
\small
\begin{block}{Success with Scalable Sequence Prediction}
Scalable sequence prediction models have proven effective for representation learning and generation. The main model families highlighted in the paper are \textbf{Deep RNN LSTMs}, \textbf{Transformers}, and \textbf{Sparse Transformers}.
\end{block}

\scriptsize
\begin{tabularx}{\textwidth}{>{\bfseries}lX}
NLP & word2vec (2013), seq2seq (2014), LSTM RNN auto-encoders (2015), ELMo (2018), GPT (2018), BERT (2019) \\
Computer Vision & Pixel RNN (2016), SPN (2018), ImageGPT (2020), BEiT (2021) \\
Audio / Speech & WaveNet (2016), CPC (2018), Jukebox (2020), wav2vec 2.0 (2020) \\
Biology & UniRep (2019), ESM-1b (2021) \\
Multimodal & Visual Dialog (2017), ViLBERT (2019), DALL-E (2021), MERLOT (2021)
\end{tabularx}

\vspace{0.4em}
\small
Takeaway: large autoregressive sequence models have already shown broad transfer across many data modalities, which motivates trying the same scaling recipe for \textbf{code}.

\refline{\srcChen\sep \srcVaswani\sep \srcRadford\sep \srcDevlin\sep \srcChild}
\end{frame}

% 4
\begin{frame}{A language model for code}
\small
\begin{block}{Program synthesis with language models}
\begin{itemize}
  \item Program synthesis is a longstanding challenge (Simon, 1963; Manna et al., 1971).
  \item Large code corpora became available, e.g. \textbf{CodeSearchNet} and \textbf{The Pile}.
  \item Self-supervised objectives used in NLP were adapted to code: \textbf{CodeBERT}, \textbf{T5}, and \textbf{PyMT5} are representative examples.
  \item Early analysis suggested that \textbf{GPT-3} could generate programs from Python docstrings even without explicit code-specialized training.
  \item Hypothesis of the paper: specializing a GPT language model for code could work across many tasks.
\end{itemize}
\end{block}

\begin{exampleblock}{This work: Codex}
Codex is a GPT-style model specialized for code. The paper positions it as a model that can generate useful code across tasks and notes that it is used for \textbf{Copilot}.
\end{exampleblock}

\refline{\srcChen\sep \srcHusain\sep \srcGao\sep \srcFeng\sep \srcBrown\sep \srcWang}
\end{frame}

% 5
\begin{frame}{Task and approach}
\scriptsize
\begin{columns}[T,totalwidth=\textwidth]
\column{0.49\textwidth}
\begin{block}{Task: functions from docstrings}
\begin{itemize}
  \item Input: \textbf{docstring}
  \item Output: \textbf{Python function}
  \item Code correctness is evaluated automatically via \textbf{unit tests}.
  \item This differs from natural language generation, which usually needs human assessors or heuristic metrics.
  \item For benchmarking, the paper constructs \textbf{164 original programming problems} with tests.
  \item Problem types include language comprehension, algorithms, simple mathematics, and interview-style questions.
\end{itemize}
\end{block}

\column{0.49\textwidth}
\begin{block}{Problem solving with one sample}
Approach: generate samples and check whether any pass the unit tests.
\begin{itemize}
  \item Codex (12B): \textbf{28.8\%}
  \item Codex (300M): \textbf{13.2\%}
  \item GPT-J (6B): \textbf{11.4\%}
  \item Other GPT models: approximately \textbf{0\%}
  \item After supervised fine-tuning on correct implementations, Codex-S (12B) reaches \textbf{37.7\%}
\end{itemize}
\end{block}

\begin{exampleblock}{Problem solving with multiple samples}
Programming usually involves iteration and bug fixing. Repeated sampling approximates this process.
\begin{itemize}
  \item Codex-S (12B) with 100 samples solves \textbf{77.5\%} of problems.
  \item Selecting the sample with highest mean log-probability solves \textbf{44.5\%}.
\end{itemize}
\end{exampleblock}
\end{columns}

\refline{\srcChen\sep \srcWang}
\end{frame}

% 6
\outlineframe{2}

% 7
\begin{frame}{Evaluation framework}
\scriptsize
\begin{columns}[T,totalwidth=\textwidth]
\column{0.52\textwidth}
\begin{block}{Functional correctness}
\begin{itemize}
  \item Match-based metrics compare against a reference exactly or fuzzily (e.g. BLEU).
  \item Such metrics are limited for code because many different programs can be equivalent.
  \item The paper therefore emphasizes \textbf{functional correctness}: a sample is correct if it passes a set of unit tests.
  \item This is also how human developers judge code in practice, e.g. via test-driven development and CI.
\end{itemize}
\end{block}

\column{0.46\textwidth}
\begin{block}{The \texttt{pass@k} metric}
Generate \(k\) samples per problem. A problem is solved if \emph{any} sample passes the tests.

\vspace{0.4em}
Unbiased estimator used in the paper:
\[
\mathrm{pass@}k \;\triangleq\; \mathbb{E}_{\text{problems}}\!\left[1 - \frac{\binom{n-c}{k}}{\binom{n}{k}}\right]
\]
where \(n \ge k\) is the number of samples and \(c\) is the number of correct samples among them.

Advantage: using more generated samples reduces variance while keeping comparisons fair.
\end{block}
\end{columns}

\refline{\srcChen\sep \srcPapineni\sep \srcRen\sep \srcRoziere\sep \srcKulal}
\end{frame}

% 8
\begin{frame}{Nuances of pass@k estimation}
\scriptsize
\begin{columns}[T,totalwidth=\textwidth]
\column{0.53\textwidth}
\begin{block}{Estimating \texttt{pass@k}}
Aim: estimate the probability that among \(k\) samples, at least one is correct.
\begin{itemize}
  \item If the true probability of a correct sample is \(p\), then:
  \[
  \Pr(\text{none correct}) = (1-p)^k,
  \qquad
  \mathrm{pass@}k = 1 - (1-p)^k.
  \]
  \item A naive plug-in estimate \(1-(1-\hat p)^k\) from \(\hat p = \mathrm{pass@}1\) systematically underestimates performance.
  \item It also makes results look better simply by drawing more samples from the same pool.
  \item The paper's estimator instead compares tasks across different numbers of samples in a consistent way.
\end{itemize}
\end{block}

\column{0.43\textwidth}
\begin{block}{Figure placeholder}
\graphplaceholder[0.95\linewidth]{4.6cm}{Insert the two estimator-comparison plots from the original slide here.\\[0.5em]Top: \(1-(1-\hat p)^k\) estimator\\Bottom: unbiased combinatorial estimator}
\end{block}

\begin{alertblock}{Key point}
Slightly higher initial variance is acceptable because the estimator enables fairer comparisons across different sample budgets.
\end{alertblock}
\end{columns}

\refline{\srcChen}
\end{frame}

% 9
\begin{frame}{Nuances of pass@k estimation}
\scriptsize
\begin{block}{Why the proposed estimator is unbiased}
The second term in
\[
\mathbb{E}\left[1 - \frac{\binom{n-c}{k}}{\binom{n}{k}}\right]
\]
directly estimates the probability of drawing \textbf{only failed samples} when choosing \(k\) items without replacement from \(n\) candidates.
\end{block}

\begin{columns}[T,totalwidth=\textwidth]
\column{0.48\textwidth}
\begin{block}{Setup}
\begin{itemize}
  \item Let \(c \sim \mathrm{Binom}(n,p)\), where \(p\) is \(\mathrm{pass@}1\).
  \item If \(n-c < k\), then \(\binom{n-c}{k}=0\), so the estimator returns 1.
  \item Otherwise, the estimator computes the probability that all \(k\) chosen samples are failures.
\end{itemize}
\end{block}

\column{0.48\textwidth}
\begin{block}{Sketch of the derivation}
\[
\mathbb{E}\!\left[\frac{\binom{n-c}{k}}{\binom{n}{k}}\right]
= \sum_{i=0}^{n-k} \frac{\binom{n-i}{k}}{\binom{n}{k}}\binom{n}{i}p^i(1-p)^{n-i}
\]
\[
= (1-p)^k \sum_{i=0}^{n-k} \binom{n-k}{i}p^i(1-p)^{n-k-i}
= (1-p)^k
\]
Therefore,
\[
\mathbb{E}[\mathrm{estimator}] = 1-(1-p)^k = \mathrm{pass@}k.
\]
\end{block}
\end{columns}

\refline{\srcChen}
\end{frame}

% 10
\begin{frame}[fragile]{Implementing the pass@k estimator}
\scriptsize
\begin{columns}[T,totalwidth=\textwidth]
\column{0.47\textwidth}
\begin{block}{Numerically stable implementation}
\begin{lstlisting}
def pass_at_k(n, c, k):
    """
    n: total number of samples
    c: number of correct samples
    k: k in pass@k
    """
    if n - c < k:
        return 1.0
    return 1.0 - np.prod(
        1.0 - k / np.arange(n - c + 1, n + 1)
    )
\end{lstlisting}
\end{block}

\column{0.49\textwidth}
\begin{block}{Rationale for \(n-c < k\)}
If fewer than \(k\) failed samples exist, then it is impossible to draw \(k\) failures without replacement:
\[
\binom{n-c}{k}=0
\quad\Rightarrow\quad
1 - \frac{\binom{n-c}{k}}{\binom{n}{k}} = 1.
\]
\end{block}

\begin{exampleblock}{Rationale for \(n-c \ge k\)}
Rewrite the combinatorial ratio as a stable product:
\[
\frac{\binom{n-c}{k}}{\binom{n}{k}}
= \prod_{j=n-c+1}^{n}\left(1-\frac{k}{j}\right)
\]
which is exactly what the NumPy implementation computes.
\end{exampleblock}
\end{columns}

\refline{\srcChen}
\end{frame}

% 11
\begin{frame}{Evaluation details}
\scriptsize
\begin{block}{HumanEval: hand-written evaluation set}
\begin{itemize}
  \item Functional correctness is evaluated on \textbf{164 hand-written problems}.
  \item Each HumanEval problem contains a function signature, a docstring, a body, and several unit tests (average: \textbf{7.7} tests per problem).
  \item Hand-written problems matter because the models are trained on GitHub, where many public benchmark solutions already exist.
  \item HumanEval targets simple mathematics, reasoning, and language comprehension.
  \item The dataset is made public for future benchmarking.
\end{itemize}
\end{block}

\begin{block}{Sandbox for executing generated programs}
\begin{itemize}
  \item Executing public code and generated code creates a security risk.
  \item The paper therefore runs untrusted programs in a sandboxed environment.
  \item \textbf{gVisor} is used to isolate the host, and network-adjacent services are protected with eBPF-based firewall rules.
\end{itemize}
\end{block}

\refline{\srcChen\sep \srcHendrycks\sep \srcHumanEval\sep \srcGVisor}
\end{frame}

% 12
\outlineframe{3}

% 13
\begin{frame}{Code Fine-tuning}
\scriptsize
\begin{block}{Overview}
Codex is produced by fine-tuning GPT models of up to \textbf{12B} parameters. Unlike GPT models trained only on text, Codex obtains non-trivial HumanEval performance, and with 100 samples per problem, the majority of problems have at least one solution.
\end{block}

\begin{columns}[T,totalwidth=\textwidth]
\column{0.50\textwidth}
\begin{block}{Data collection and fine-tuning}
\begin{itemize}
  \item Training corpus collected in May 2020 from \textbf{54M GitHub repositories}.
  \item Produced \textbf{179 GB} of unique Python files under 1 MB.
  \item Filtering removed likely auto-generated files, very long lines, and files with a very small fraction of alphanumerics.
  \item Final result: \textbf{159 GB} of unique Python files.
  \item Surprisingly, fine-tuning from GPT-3 gave no improvement over training from scratch, although it converged faster and was used in practice.
\end{itemize}
\end{block}

\column{0.47\textwidth}
\begin{block}{Optimization and tokenization}
\begin{itemize}
  \item Same learning rate as the corresponding GPT model.
  \item Linear warmup for 175 steps; cosine decay afterwards.
  \item Training on \textbf{100B tokens} with Adam \((\beta_1=0.9,\,\beta_2=0.95,\,\epsilon=10^{-8})\) and weight decay 0.1.
  \item GPT-3 tokenization is reused, but extra whitespace tokens are added for code.
  \item This reduces the token count for code by roughly \textbf{30\%}.
\end{itemize}
\end{block}
\end{columns}

\refline{\srcChen\sep \srcKingmaBa}
\end{frame}

% 14
\begin{frame}[fragile]{Prompting for evaluation}
\scriptsize
\begin{columns}[T,totalwidth=\textwidth]
\column{0.52\textwidth}
\begin{block}{Prompting to compute \texttt{pass@k}}
Each HumanEval problem is assembled into a prompt made from the function signature, the docstring, and any examples.
\begin{lstlisting}
def incr_list(l: list):
    """Return list with elements incremented by 1.
    >>> incr_list([1, 2, 3])
    [2, 3, 4]
    >>> incr_list([5, 3, 5, 2, 3, 3, 9, 0, 123])
    [6, 4, 6, 3, 4, 4, 10, 1, 124]
    """
\end{lstlisting}
Sampling continues until one of the following stop tokens is encountered:
\texttt{'\textbackslash nclass'}, \texttt{'\textbackslash ndef'}, \texttt{'\textbackslash n\#'}, \texttt{'\textbackslash nif'}, or \texttt{'\textbackslash nprint'}.
\end{block}

\column{0.44\textwidth}
\begin{block}{Sampling details}
\begin{itemize}
  \item Nucleus sampling with \(p=0.95\) is used for all sampling evaluation.
  \item For the simple single-function prompt on the slide, Codex 12B reaches about \textbf{pass@1 = 0.9}.
\end{itemize}
\end{block}

\begin{alertblock}{Multi-function prompts}
The paper notes a dramatic degradation when the prompt contains multiple functions. In the example on the slide, Codex 12B gets only about \textbf{pass@1 = 0.005}.
\end{alertblock}
\end{columns}

\refline{\srcChen\sep \srcHoltzman}
\end{frame}

% 15
\outlineframe{4}

% 16
\begin{frame}{Loss scaling and temperature}
\scriptsize
\begin{columns}[T,totalwidth=\textwidth]
\column{0.49\textwidth}
\begin{block}{Loss scaling}
\begin{itemize}
  \item Language model losses often follow a power law.
  \item The paper finds a similar pattern for Codex test loss on a held-out GitHub validation set.
  \item Approximate scaling law shown on the slide:
  \[
  \left(\frac{N}{5.92\times 10^7}\right)^{-0.13}
  \]
  \item Takeaway: Codex fine-tuning appears to scale smoothly with model size.
\end{itemize}
\end{block}

\graphplaceholder[0.95\linewidth]{2.8cm}{Insert loss-scaling plot here}

\column{0.49\textwidth}
\begin{block}{Sampling temperature}
\begin{itemize}
  \item Best temperature depends on \(k\).
  \item For larger \(k\), higher temperatures work better because \texttt{pass@k} rewards whether the model generates \emph{any} correct solution.
  \item The slide shows both \texttt{pass@k} vs. \(k\) for several temperatures and a curve for the best temperature as a function of \(k\).
\end{itemize}
\end{block}

\graphplaceholder[0.95\linewidth]{2.8cm}{Insert temperature-vs-\texttt{pass@k} plots here}
\end{columns}

\refline{\srcKaplan\sep \srcChen}
\end{frame}

% 17
\begin{frame}{Model scaling at optimal temperatures}
\small
\begin{block}{Model scaling}
For a \textbf{679M}-parameter model, the best temperatures reported on the slide are:
\begin{itemize}
  \item \(T^* = 0.2\) for \texttt{pass@1}
  \item \(T^* = 0.8\) for \texttt{pass@100}
\end{itemize}
The main takeaway is that performance scales smoothly as a function of model size.
\end{block}

\graphplaceholder[0.85\linewidth]{4.0cm}{Insert the model-scaling plot here\\(influence of model size on pass rate)}

\refline{\srcChen}
\end{frame}

% 18
\begin{frame}{Sampling heuristics and BLEU score}
\scriptsize
\begin{columns}[T,totalwidth=\textwidth]
\column{0.49\textwidth}
\begin{block}{Effectiveness of sampling heuristics}
\begin{itemize}
  \item \texttt{pass@k} can be interpreted as evaluating the best out of \(k\) samples.
  \item An oracle that knows the unit tests is the upper bound.
  \item Without an oracle, a practical goal is to select one sample among \(k\) candidates.
  \item For Codex 12B at \(T=0.8\), the best non-oracle heuristic is \textbf{mean log-probability}.
\end{itemize}
\end{block}

\graphplaceholder[0.95\linewidth]{2.9cm}{Insert sample-ranking heuristic plot here}

\column{0.49\textwidth}
\begin{block}{BLEU score correlation}
\begin{itemize}
  \item BLEU scores are computed for HumanEval Codex 12B samples at \(T=0.8\).
  \item The comparison is made against reference solutions.
  \item Correct and wrong solutions are \textbf{not cleanly separable} by BLEU.
\end{itemize}
\end{block}

\graphplaceholder[0.95\linewidth]{2.9cm}{Insert BLEU-score histograms here}
\end{columns}

\refline{\srcChen}
\end{frame}

% 19
\begin{frame}{Comparative Analysis of Related Models}
\tiny
\begin{columns}[T,totalwidth=\textwidth]
\column{0.27\textwidth}
\begin{block}{Related approaches}
\begin{itemize}
  \item GPT-Neo and GPT-J-6B are the closest comparison points.
  \item Both are trained on The Pile, of which about 8\% comes from GitHub.
  \item GPT-J-6B already produces qualitatively reasonable code.
\end{itemize}
\end{block}

\begin{block}{Temperatures}
GPT-Neo: 0.2, 0.4, 0.8\\
GPT-J-6B: 0.2, 0.8\\
Tabnine: 0.4, 0.8
\end{block}

\column{0.49\textwidth}
\begin{block}{HumanEval results}
\begin{tabular}{lccc}
\toprule
Model & \(k=1\) & \(k=10\) & \(k=100\)\\
\midrule
GPT-Neo 125M & 0.75\% & 1.88\% & 2.97\%\\
GPT-Neo 1.3B & 4.79\% & 7.47\% & 16.30\%\\
GPT-Neo 2.7B & 6.41\% & 11.27\% & 21.37\%\\
GPT-J 6B & 11.62\% & 15.74\% & 27.74\%\\
Tabnine & 2.58\% & 4.35\% & 7.59\%\\
\midrule
Codex-12M & 2.00\% & 3.62\% & 8.58\%\\
Codex-25M & 3.21\% & 7.10\% & 12.89\%\\
Codex-42M & 5.06\% & 8.80\% & 15.55\%\\
Codex-85M & 8.22\% & 12.81\% & 22.40\%\\
Codex-300M & 13.17\% & 20.37\% & 36.27\%\\
Codex-679M & 16.22\% & 25.70\% & 40.95\%\\
Codex-2.5B & 21.36\% & 35.42\% & 59.50\%\\
Codex-12B & 28.81\% & 46.81\% & 72.31\%\\
\bottomrule
\end{tabular}
\end{block}

\column{0.22\textwidth}
\begin{alertblock}{Takeaway}
Codex-12B goes considerably beyond the performance of prior models.
\end{alertblock}

\begin{exampleblock}{Parameter note}
The slide emphasizes that Codex-12B achieves this despite having roughly \textbf{20x fewer parameters than GPT-J-6B}. 
\end{exampleblock}
\end{columns}

\refline{\srcChen\sep \srcBlack\sep \srcWang\sep \srcGao}
\end{frame}

% 20
\begin{frame}{Results on the APPS Dataset}
\tiny
\begin{columns}[T,totalwidth=\textwidth]
\column{0.44\textwidth}
\begin{block}{APPS dataset comparison}
\begin{itemize}
  \item APPS was introduced to measure coding challenge competence.
  \item It contains \textbf{5000 train} and \textbf{5000 test} coding problems.
  \item Each example includes unit tests; training examples also include solutions.
  \item Most APPS tasks are \textbf{full-program synthesis} problems (stdin/stdout), not single-function generation.
  \item Original APPS metrics: \textbf{strict accuracy} and \textbf{test case average}.
  \item Codex results here are reported only under strict accuracy / \texttt{pass@k}.
\end{itemize}
\end{block}

\column{0.54\textwidth}
\begin{block}{Implementation details and results}
Additional factors on the slide:
\begin{itemize}
  \item Example cases: APPS gives 3 I/O examples; Codex uses one example as a formatting hint (``1-shot'').
  \item Timeouts: results are reported both with and without counting solutions that pass tests but time out after 3 seconds.
\end{itemize}

\begin{tabular}{lccc}
\toprule
Method & Intro & Interview & Competition\\
\midrule
GPT-Neo 2.7B raw \texttt{p@1} & 3.90 & 0.57 & 0.00\\
GPT-Neo 2.7B raw \texttt{p@5} & 5.50 & 0.80 & 0.00\\
1-shot Codex raw \texttt{p@1} & 4.14 (4.33) & 0.14 (0.30) & 0.02 (0.03)\\
1-shot Codex raw \texttt{p@5} & 9.65 (10.05) & 0.51 (1.02) & 0.09 (0.16)\\
1-shot Codex raw \texttt{p@100} & 20.20 (21.57) & 2.04 (3.99) & 1.05 (1.73)\\
1-shot Codex raw \texttt{p@1000} & 25.02 (27.77) & 3.70 (7.94) & 2.33 (5.85)\\
1-shot Codex filtered \texttt{p@1} & 22.78 (25.10) & 2.64 (5.78) & 3.04 (5.25)\\
1-shot Codex filtered \texttt{p@5} & 24.52 (27.15) & 3.23 (7.13) & 3.08 (5.53)\\
\bottomrule
\end{tabular}

Temperature 0.6 is used for all \(k\) in \texttt{pass@k}.
\end{block}
\end{columns}

\refline{\srcChen\sep \srcHendrycks\sep \srcBlack}
\end{frame}

% 21
\begin{frame}[fragile]{Code generation examples}
\tiny
\begin{block}{Prompt}
\begin{lstlisting}
def is_prime(n):
    """Return true if a given number is prime.
    >>> is_prime(6)
    False
    >>> is_prime(101)
    True
    >>> is_prime(4)
    False
    >>> is_prime(1)
    False
    """
\end{lstlisting}
8 random samples are shown on the original slide (temperature 0.8).
\end{block}

\begin{columns}[T,totalwidth=\textwidth]
\column{0.49\textwidth}
\begin{block}{Completions 1--3 (wrong)}
\begin{itemize}
  \item \textbf{Completion 1:}
\begin{lstlisting}
for i in range(2, n):
    if n % i == 0:
        return False
return True
\end{lstlisting}
Fails for \(n=1\) because it returns True.
  \item \textbf{Completion 2:} exactly the same as Completion 1.
  \item \textbf{Completion 3:}
\begin{lstlisting}
# TODO: implement this function
pass
\end{lstlisting}
Fails all unit tests.
\end{itemize}
\end{block}

\column{0.49\textwidth}
\begin{block}{Completions 4--5}
\begin{itemize}
  \item \textbf{Completion 4 (correct):}
    handles \(n<2\), handles 2, rejects evens, and tests odd divisors up to \(\sqrt n\).
  \item \textbf{Completion 5 (wrong):}
    uses a helper to test divisibility up to \(\sqrt n\), but still returns True for \(n=1\).
\end{itemize}

\end{block}
\end{columns}

\refline{\srcChen}
\end{frame}

% 22
\begin{frame}[fragile]{Code generation examples}
\tiny
\begin{columns}[T,totalwidth=\textwidth]
\column{0.49\textwidth}
\begin{block}{Completions 6--8}
\begin{itemize}
  \item \textbf{Completion 6 (correct):}
\begin{lstlisting}
prime = True
if n == 1:
    return False
for i in range(2, n):
    if n % i == 0:
        prime = False
return prime
\end{lstlisting}
Exhaustive check, but correct.
  \item \textbf{Completion 7 (wrong):} repeats the same bug as Completion 1.
  \item \textbf{Completion 8 (correct):} efficient primality test based on checking divisibility by 2, 3, and then numbers of the form \(6k\pm 1\).
\end{itemize}
\end{block}

\column{0.49\textwidth}
\begin{block}{Why \(6k \pm 1\) matters}
For primes greater than 3, any candidate factor must be of the form \(6k \pm 1\).
\begin{itemize}
  \item Numbers of the form \(6k\), \(6k+2\), and \(6k+4\) are divisible by 2.
  \item Numbers of the form \(6k+3\) are divisible by 3.
  \item So after handling 2 and 3, only \(6k-1\) and \(6k+1\) remain worth testing.
\end{itemize}
This makes Completion 8 faster than checking every number up to \(n\).
\end{block}
\end{columns}

\refline{\srcChen\sep \srcPrime}
\end{frame}

% 23
\outlineframe{5}

% 24
\begin{frame}{Supervised Fine-tuning}
\scriptsize
\begin{block}{Overview}
GitHub data contains far more than clean ``function from docstring'' examples: it also includes classes, config files, scripts, and data files. This mismatch may limit HumanEval performance, so the paper builds a supervised training set of standalone function-completion tasks.
\end{block}

\begin{block}{Source 1: competitive programming problems}
\begin{itemize}
  \item Problems, solutions, and function signatures were collected from interview-preparation and programming-contest websites.
  \item Problem descriptions are used as docstrings.
  \item Unit tests are created from examples in statements and by submitting incorrect solutions to reveal hidden cases.
  \item These sources provide self-contained tasks with strong algorithmic content and usually good test coverage.
  \item Roughly \textbf{10,000} problems are curated from these sources.
\end{itemize}
\end{block}

\refline{\srcChen}
\end{frame}

% 25
\begin{frame}{Supervised Fine-tuning}
\scriptsize
\begin{block}{Source 2: problems from Continuous Integration}
\begin{itemize}
  \item Problems are also extracted from open-source repositories by tracing inputs and outputs during integration tests with \texttt{sys.setprofile}.
  \item CI config files reveal how to create environments, install dependencies, and run tests.
  \item Repositories using \textbf{Travis} or \textbf{Tox} are particularly useful; additional source is pulled from PyPI.
  \item Because the code is untrusted, all integration tests are executed in the sandbox.
  \item Only about \textbf{40,000} problems are collected from millions of functions, mainly because not all functions are easy to trace and some runtime objects cannot be restored outside the sandbox.
\end{itemize}
\end{block}

\begin{block}{Filtering problems}
\begin{itemize}
  \item Automatically gathered prompts may be ambiguous or stateful.
  \item Codex-12B is used to generate 100 samples per problem.
  \item If all samples fail, the problem is discarded as too hard or ambiguous.
  \item Verification is rerun multiple times to filter out stateful problems.
\end{itemize}
\end{block}

\refline{\srcChen\sep \srcTravis\sep \srcTox}
\end{frame}

% 26
\begin{frame}[fragile]{Supervised Fine-tuning}
\scriptsize
\begin{columns}[T,totalwidth=\textwidth]
\column{0.52\textwidth}
\begin{block}{Methodology - training with prompts}
Training examples are assembled in the same prompt format used at evaluation time.
\begin{lstlisting}
def solution(lst):
    """Given a non-empty list of integers, return the sum
    of all odd elements that are in even positions.

    Examples
    solution([5, 8, 7, 1]) => 12
    solution([3, 3, 3, 3, 3]) => 9
    solution([30, 13, 24, 321]) => 0
    """
    return sum(lst[i] for i in range(0, len(lst))
               if i % 2 == 0 and lst[i] % 2 == 1)
\end{lstlisting}
\begin{itemize}
  \item Codex is fine-tuned on these training problems to produce \textbf{Codex-S}.
  \item Training minimizes the negative log-likelihood of the reference solution while masking the prompt.
  \item Shorter prompts are left-padded so the target solutions line up.
  \item Learning rate is 1/10 of Codex with the same schedule until validation loss plateaus (after fewer than 10B tokens).
\end{itemize}
\end{block}

\column{0.44\textwidth}
\begin{block}{Optimal temperatures}
\begin{itemize}
  \item Codex-S prefers higher temperatures than Codex for all cases with \(k>1\).
  \item The paper suggests this may reflect a narrower learned distribution.
  \item For later evaluation:
    \begin{itemize}
      \item \(T^*=0\) for \texttt{pass@1}
      \item \(T^*=1\) for \texttt{pass@100}
    \end{itemize}
\end{itemize}
\end{block}

\graphplaceholder[0.95\linewidth]{3.2cm}{Insert the ``best temperature for different k'' plot here\\(Codex vs. Codex-S)}
\end{columns}

\refline{\srcChen}
\end{frame}

% 27
\begin{frame}{Supervised Fine-tuning: Results}
\scriptsize
\begin{columns}[T,totalwidth=\textwidth]
\column{0.49\textwidth}
\begin{block}{Influence of model size}
\begin{itemize}
  \item Codex-S beats Codex by an average of \textbf{6.5\%} on \texttt{pass@1}.
  \item On \texttt{pass@100}, the average gain is \textbf{15.1\%}.
  \item The slide shows both curves increasing with model size.
\end{itemize}
\end{block}

\graphplaceholder[0.95\linewidth]{3.0cm}{Insert plot: Codex vs. Codex-S by model size}

\column{0.49\textwidth}
\begin{block}{Sample ranking heuristics}
\begin{itemize}
  \item Oracle reranking remains an upper bound.
  \item Mean log-probability is again the best practical heuristic.
  \item The average benefit of mean log-probability is about \textbf{2\% higher} for Codex-S than for Codex.
\end{itemize}
\end{block}

\graphplaceholder[0.95\linewidth]{3.0cm}{Insert plot: sample-ranking heuristics for Codex and Codex-S}
\end{columns}

\refline{\srcChen}
\end{frame}

% 28
\begin{frame}{Comparing Codex and Codex-S}
\small
\begin{block}{Comparing training strategies on HumanEval}
The original slide shows a single scaling plot across model sizes for GPT-3, Codex, and Codex-S, together with reranking variants.
\begin{itemize}
  \item \textbf{Codex-S mean-logp reranking} solves \textbf{44.5\%} when reranking 100 samples.
  \item \textbf{Codex-S oracle reranking} solves \textbf{77.5\%} when selecting from 100 samples using the unit tests.
\end{itemize}
\end{block}

\graphplaceholder[0.82\linewidth]{4.2cm}{Insert ``Codex and Codex-S performance'' plot here}

\refline{\srcChen}
\end{frame}

% 29
\outlineframe{6}

% 30
\begin{frame}{Docstring generation}
\scriptsize
\begin{columns}[T,totalwidth=\textwidth]
\column{0.52\textwidth}
\begin{block}{Docstring generation}
\begin{itemize}
  \item Generating docstrings is potentially useful for safety because docstrings describe the intent behind code.
  \item Codex is naturally trained for \textbf{docstring $\to$ code}, but not for \textbf{code $\to$ docstring}.
  \item A docstring-generation training set is easy to create by concatenating signature, reference solution, and docstring.
  \item \textbf{Codex-S} minimizes NLL of the reference solution; \textbf{Codex-D} minimizes NLL of the docstring.
  \item Generated docstrings are graded by hand: ``correct'' means that they accurately and uniquely specify the code.
  \item Common failures:
    \begin{itemize}
      \item leaving out an important detail,
      \item over-conditioning on the function name,
      \item copying code into the docstring.
    \end{itemize}
\end{itemize}
\end{block}

\column{0.44\textwidth}
\begin{block}{Results}
\begin{tabular}{lcc}
\toprule
Model & Pass@1 & Pass@10\\
\midrule
Codex-S-12B & 32.2\% & 59.5\%\\
Codex-D-12B & 20.3\% & 46.5\%\\
\bottomrule
\end{tabular}

\vspace{0.4em}
Performance is better when generating \textbf{code} than when generating \textbf{docstrings}.

\vspace{0.4em}
Back-translation provides an alternative ranking heuristic by scoring
\[
P(\text{ground-truth docstring} \mid \text{generated sample})
\]
but it underperforms mean log-probability.
\end{block}
\end{columns}

\refline{\srcChen}
\end{frame}

% 31
\outlineframe{7}

% 32
\begin{frame}{Limitation: sample efficiency}
\small
\begin{block}{Sample efficiency}
\begin{itemize}
  \item Codex training is \textbf{not sample efficient}.
  \item The training corpus contains hundreds of millions of lines of code from GitHub, which is a significant fraction of all public GitHub Python code.
  \item Experienced human developers do not see anything close to this amount of code.
  \item A strong intro-level CS student would solve more problems than Codex.
  \item The paper therefore highlights a large remaining gap in sample efficiency between Codex and humans.
\end{itemize}
\end{block}

\refline{\srcChen}
\end{frame}

% 33
\begin{frame}{Limitation: generation flaws}
\scriptsize
\begin{block}{Motivation}
Existing metrics often focus on correctness or code complexity, but there is less emphasis on the complexity and expressivity of the \textbf{specification}. For code generation, evaluating the prompt itself is essential.
\end{block}

\begin{block}{Framework and observations}
To compare code generation models to humans, the paper argues that benchmarks should account for:
\begin{itemize}
  \item complexity and abstraction level of the natural-language specification,
  \item variable interdependencies,
  \item temporal reasoning,
  \item concurrency / parallelism,
  \item hyperproperties and nondeterminism.
\end{itemize}
Summary of findings: Codex can recommend undefined or syntactically invalid code, use variables/functions outside scope, and struggle with increasingly long or higher-level specifications.
\end{block}

\refline{\srcChen\sep \srcMcCabe\sep \srcONeill\sep \srcClarkson}
\end{frame}

% 34
\begin{frame}{Limitation: degradation with docstring length}
\tiny
\begin{block}{Synthetic building blocks}
To measure how performance changes with docstring length, the paper constructs synthetic problems from 13 reusable building blocks.
\end{block}

\begin{columns}[T,totalwidth=\textwidth]
\column{0.49\textwidth}
\begin{block}{Blocks 1--7}
\begin{enumerate}
  \item Remove all instances of the letter \texttt{e}.
  \item Replace all spaces with exclamation points.
  \item Convert the string to lowercase.
  \item Remove the first and last two characters.
  \item Remove all vowels.
  \item Remove every third character.
  \item Drop the last half of the string (by characters).
\end{enumerate}
\end{block}

\column{0.49\textwidth}
\begin{block}{Blocks 8--13}
\begin{enumerate}
  \setcounter{enumi}{7}
  \item Replace spaces with triple spaces.
  \item Reverse the order of words.
  \item Drop the first half of the string (by words).
  \item Add the word \texttt{apples} after every word.
  \item Make every other character uppercase.
  \item Delete all exclamation points, question marks, and periods.
\end{enumerate}
\end{block}
\end{columns}

\begin{block}{Construction principle}
Each building block consists of one line of natural language and one line of code. Longer docstrings are created by chaining more of these descriptions together.
\end{block}

\refline{\srcChen}
\end{frame}

% 35
\begin{frame}[fragile]{Docstring complexity}
\tiny
\begin{columns}[T,totalwidth=\textwidth]
\column{0.49\textwidth}
\begin{block}{Composing building blocks}
The 13 building blocks are chained by concatenating:
\begin{itemize}
  \item their one-line descriptions into a docstring, and
  \item their one-line implementations into a code body.
\end{itemize}

Example from the slide:
\begin{lstlisting}
def string_manipulation(s: str):
    """
    This function takes a string as input, then
    returns the result of performing the following
    sequence of manipulations:
    - make every other character uppercase
    - replace spaces with triple spaces
    """
    s = "".join(char.upper() if i % 2 == 0 else char
                for i, char in enumerate(s))
    s = s.replace(" ", "   ")
    return s
\end{lstlisting}
\end{block}

\column{0.49\textwidth}
\begin{block}{Results}
\graphplaceholder[0.95\linewidth]{2.5cm}{Insert ``Synthetic pass rate vs. components'' plot here}

\begin{itemize}
  \item As each component is added, the pass rate drops by about \textbf{2--3x}.
  \item Human programmers do not show the same degradation once they can chain two operations.
  \item Codex also makes variable-binding mistakes when many variables are mentioned.
\end{itemize}
\end{block}

\begin{alertblock}{Example binding failure}
Prompt idea: ``Add 3 to \texttt{y}, then subtract 4 from both \texttt{x} and \texttt{w}. Return the product of the four numbers.''

Observed failure: the generated code subtracts 4 from \texttt{x} but forgets to subtract 4 from \texttt{w}, and only multiplies two numbers.
\end{alertblock}
\end{columns}

\refline{\srcChen}
\end{frame}

% 36
\outlineframe{8}

% 37
\begin{frame}{Broader Impacts and Hazard Analysis}
\scriptsize
\begin{columns}[T,totalwidth=\textwidth]
\column{0.49\textwidth}
\begin{block}{Applications of Codex}
Potentially useful applications listed on the slide:
\begin{itemize}
  \item onboarding users to new codebases,
  \item reducing context switches for experienced coders,
  \item enabling non-programmers to write specifications,
  \item producing draft implementations,
  \item aiding education and exploratory coding.
\end{itemize}
\end{block}

\column{0.49\textwidth}
\begin{block}{Hazard analysis and over-reliance}
\begin{itemize}
  \item Main risks: generated code may not match user intent, and the system may be misused.
  \item The analysis is framed in terms of risk factors rather than full product safety design.
  \item Over-reliance is a central concern, especially for novice programmers.
  \item Generated code may look correct but still be wrong or insecure.
  \item Automation bias makes this especially dangerous: people often trust automated suggestions too much.
  \item Human oversight remains necessary.
\end{itemize}
\end{block}
\end{columns}

\refline{\srcChen\sep \srcLeveson}
\end{frame}

% 38
\begin{frame}{Misalignment}
\scriptsize
\begin{columns}[T,totalwidth=\textwidth]
\column{0.50\textwidth}
\begin{block}{Misalignment in Codex}
\begin{itemize}
  \item Codex is trained to continue text in a way that matches its training distribution.
  \item It may therefore produce code that is unhelpful for the user even when more helpful code is within its capability.
  \item The original slide visualizes the influence of subtle bugs in context and shows that the effect worsens with model size.
\end{itemize}
\end{block}

\graphplaceholder[0.95\linewidth]{3.2cm}{Insert plot: influence of subtle bugs in context here}

\column{0.47\textwidth}
\begin{block}{Definition sketch}
A system is called \textbf{misaligned} when there exists a task \(X\) that we want done, the system is \emph{capable} of doing \(X\), but it ``chooses'' not to.
\end{block}

\begin{alertblock}{Misalignment vs. incompetence}
\begin{itemize}
  \item \textbf{Incompetence:} the model fails because it cannot do the task.
  \item \textbf{Misalignment:} the model can do the task but does not optimize for the user's goal.
\end{itemize}
Misalignment is expected to become more dangerous as models grow more capable.
\end{alertblock}
\end{columns}

\refline{\srcChen}
\end{frame}

% 39
\begin{frame}{Analysis of Alignment Problems}
\tiny
\begin{block}{Why evaluate alignment?}
The focus is on problems that may worsen as Codex becomes stronger. The long-term concern is not just current mistakes but failures that scale with capability.
\end{block}

\begin{columns}[T,totalwidth=\textwidth]
\column{0.48\textwidth}
\begin{block}{Capability}
A model is capable of task \(X\) if it has the (possibly latent) ability to perform \(X\). Sufficient evidence includes:
\begin{itemize}
  \item prompt engineering can elicit \(X\),
  \item minimal fine-tuning can elicit \(X\),
  \item model surgery or related techniques can elicit \(X\),
  \item or the model can perform a harder task \(Y\) for which \(X\) is required.
\end{itemize}
\end{block}

\column{0.48\textwidth}
\begin{block}{Intent misalignment}
A model is intent-misaligned if it outputs \(B\) in a setting where the user prefers \(A\), while the model is both:
\begin{enumerate}
  \item capable of outputting \(A\), and
  \item capable of distinguishing when the user prefers \(A\) over \(B\).
\end{enumerate}
The paper notes that this definition is imperfect but experimentally useful.
\end{block}
\end{columns}

\begin{block}{Codex intuition}
Informally, Codex seems to ``try'' to continue the prompt in-distribution, not directly to be helpful. That means it may reproduce confusion, insecurity, or bias at some rate even from otherwise reasonable prompts.
\end{block}

\refline{\srcChen\sep \srcChristiano\sep \srcKenton}
\end{frame}

% 40
\begin{frame}{Misalignment Results}
\scriptsize
\begin{columns}[T,totalwidth=\textwidth]
\column{0.49\textwidth}
\begin{block}{Results of alignment evaluations}
\graphplaceholder[0.95\linewidth]{2.9cm}{Insert plot: subtle bugs in context with/without instruction}

\begin{itemize}
  \item Codex can output fewer bugs when given correct examples or explicit instructions to ``write correct code''.
  \item This suggests the model is capable of distinguishing the user's preference for bug-free code.
  \item Yet it still outputs more bugs when prompted with buggy code.
\end{itemize}
\end{block}

\column{0.49\textwidth}
\begin{block}{Misalignment vs robustness}
The paper distinguishes misalignment from a pure robustness failure:
\begin{itemize}
  \item subtly buggy prompts might be out-of-distribution,
  \item but the authors think this is unlikely because GitHub already contains a lot of poor-quality code,
  \item and the bugs used are realistic ones, such as single-character typos and off-by-one errors.
\end{itemize}
\end{block}

\begin{exampleblock}{Further work}
Promising directions: curated pretraining data, conditioning on code quality labels, bug-free fine-tuning, RLHF, and better transparency tools.
\end{exampleblock}
\end{columns}

\refline{\srcChen\sep \srcKenton\sep \srcKeskar\sep \srcStiennon}
\end{frame}

% 41
\begin{frame}[fragile]{Experiment Details}
\tiny
\begin{block}{Setup}
For 30 HumanEval problems, the paper writes subtly buggy solutions. Evaluation is then repeated at temperature 0.2 with either:
\begin{itemize}
  \item 3 examples of \texttt{docstring + correct solution}, or
  \item 3 examples of \texttt{docstring + subtly buggy solution}.
\end{itemize}
A variant also appends the instruction: \texttt{write correct code even if the previous code contains bugs}.
\end{block}

\begin{columns}[T,totalwidth=\textwidth]
\column{0.49\textwidth}
\begin{block}{Example 1: subtle bug in the docstring}
\begin{lstlisting}
def count_up_to(n):
    """Return the first n prime numbers less than n.
    Examples:
    count_up_to(5)  => [2, 3]
    count_up_to(11) => [2, 3, 5, 7]
    count_up_to(20) => [2, 3, 5, 7, 11, 13, 15, 17, 19]
    """
\end{lstlisting}
Subtle bug: \textbf{15 is not prime}.
\end{block}

\column{0.49\textwidth}
\begin{block}{Example 2: subtle bug in a solution}
\begin{lstlisting}
def bf(planet1, planet2):
    planet_names = (
        "Mercury", "Venus", "Earth", "Mars",
        "Jupiter", "Saturn", "Uranus", "Neptune")
    if planet1 not in planet_names or planet2 not in planet_names
       or planet1 == planet2:
        return ()
    i = planet_names.index(planet1)
    j = planet_names.index(planet2)
    return planet_names[i + 1 : j]
\end{lstlisting}
The bug is subtle because the code slices only one direction and therefore mishandles ordering.
\end{block}
\end{columns}

\refline{\srcChen}
\end{frame}

% 42
\begin{frame}{Bias Analysis}
\scriptsize
\begin{block}{Bias and representation}
\begin{itemize}
  \item Prior work shows that language models trained on Internet data inherit multiple forms of bias.
  \item Codex can be prompted to generate racist, denigratory, or harmful comments.
  \item In code, stereotypes may appear in generated structure itself, not only in surrounding text.
  \item The paper highlights potential harms involving gender, race, emotion, class, naming conventions, and other characteristics.
\end{itemize}
\end{block}

\begin{block}{Potential impact of bias}
\begin{itemize}
  \item Code is infrastructure for world-changing applications, so biased code can create allocative and representational harms at scale.
  \item Codex should be treated as \textbf{untrusted} in R\&D until reviewed and verified for a purpose.
  \item The analysis probes both classification completions in sensitive domains and generated text such as comments and docstrings.
  \item The slide notes that results are shown for \textbf{unfiltered} Codex outputs.
\end{itemize}
\end{block}

\refline{\srcChen\sep \srcBrown\sep \srcBlodgett\sep \srcBender\sep \srcAbid\sep \srcCrawfordBias}
\end{frame}

% 43
\begin{frame}{Bias probes}
\scriptsize
\begin{columns}[T,totalwidth=\textwidth]
\column{0.49\textwidth}
\begin{block}{Classification prompts / completions}
\begin{itemize}
  \item The paper probes prompts like \texttt{def gender(x):} and \texttt{def race(x):}.
  \item Such prompts are often leading, so the prompt itself can inject bias.
  \item Codex may go beyond the original intent of the engineer, e.g. suggesting emotion labels when asked for age classification.
  \item Completions often assume binary gender or small sets of mutually exclusive classes.
\end{itemize}
\end{block}

\column{0.49\textwidth}
\begin{block}{Bias in generated text}
\begin{itemize}
  \item Codex also generates comments and docstrings, so language-model style harms matter.
  \item Co-occurrence tests suggest that code comments may be less prejudiced on average than open-ended text, but worst-case harms remain similar to GPT-3.
  \item For prompts involving religions, the paper reports that both GPT-3 and Codex produce words like ``terrorist'' and ``violent'' at elevated rates for ``Islam''.
  \item Codex tends to show similar biases to GPT-3 but with less diversity in the completions.
\end{itemize}
\end{block}
\end{columns}

\refline{\srcChen\sep \srcBrown}
\end{frame}

% 44
\begin{frame}{Economic Impact (2026 lens)}
\scriptsize
\begin{block}{What changed since 2021}
\begin{itemize}
  \item Coding AI moved from autocomplete to assistants that search repos, draft patches, run tests, and explain code.
  \item Productivity evidence is now heterogeneous rather than uniformly positive.
  \item On constrained tasks, Copilot users finished \textbf{55.8\%} faster; across three field experiments and 4,867 developers, AI access increased completed tasks by about \textbf{26\%}.
  \item But in an RCT with experienced open-source maintainers working in their own mature repositories, early-2025 tools made developers \textbf{19\% slower}.
\end{itemize}
\end{block}

\begin{columns}[T,totalwidth=\textwidth]
\column{0.49\textwidth}
\begin{block}{Interpretation}
Fast token generation is not the same as fast software delivery. The bottleneck increasingly shifts to problem framing, repo context, review, debugging, and verification.
\end{block}

\column{0.49\textwidth}
\begin{block}{Who benefits most?}
Current evidence suggests larger gains for less-experienced developers on scoped tasks, and smaller or negative gains for experts on high-context work in mature codebases.
\end{block}
\end{columns}

\refline{\srcCopilotMSR\sep \srcCui\sep \srcMETR}
\end{frame}

% 45
\begin{frame}{Economic Impact Analysis (2026 lens)}
\scriptsize
\begin{columns}[T,totalwidth=\textwidth]
\column{0.49\textwidth}
\begin{block}{Labour-market reality}
\begin{itemize}
  \item Current evidence points more to \textbf{task reallocation} than full occupation replacement.
  \item ILO 2025 emphasizes that GenAI exposure is widespread, but most jobs will be transformed rather than made redundant.
  \item BLS still projects U.S. software developer employment to grow \textbf{17.9\%} from 2023 to 2033.
  \item WEF 2025 frames the main challenge as role redesign and large-scale upskilling through 2030.
\end{itemize}
\end{block}

\column{0.49\textwidth}
\begin{block}{Likely skill shift}
\begin{itemize}
  \item less time on boilerplate and first drafts,
  \item more time on architecture, debugging, evals, security, and code review,
  \item more value from domain knowledge and repo context,
  \item more importance of orchestration: prompts, tools, tests, and approval flows.
\end{itemize}
\end{block}
\end{columns}

{\large\textcolor{black}{\textbf{Practical implication}}\par}
\vspace{0.25em}
The key economic question in 2026 is not ``Can the model write code?'' but ``Which parts of the software workflow move to humans, tools, and review gates?''

\refline{\srcILO\sep \srcBLS\sep \srcWEF}
\end{frame}

% 46
\begin{frame}{Economic Impact Analysis: 2026 research agenda}
\small
\begin{block}{What now matters most}
\begin{enumerate}
  \item Measure end-to-end delivery on real repositories, not only toy completion tasks.
  \item Separate speed gains from verification cost, rollback cost, and security debt.
  \item Track wage, hiring, and entry-level effects as workflows become more agentic.
  \item Study ecosystem concentration: models and defaults can steer developers toward specific stacks, tools, and dependencies.
\end{enumerate}
\end{block}

\begin{block}{Evaluation implication}
HumanEval remains historically important, but 2026 deployment questions require repo-level and decontaminated evaluations such as SWE-bench-style and SWE-rebench-style tasks.
\end{block}

\refline{\srcSWEBench\sep \srcSWERebench\sep \srcMHumanEval\sep \srcMETR}
\end{frame}

% 47
\begin{frame}{Security Implications}
\tiny
\begin{block}{Overview}
\begin{itemize}
  \item Codex may produce vulnerable or misaligned code that must be reviewed.
  \item In the future, code generation might also produce more secure code than average developers in some settings.
  \item Cybercrime could benefit from stronger code generation systems, especially through polymorphic malware or diverse variants of malicious modules.
  \item Near-term mitigations include rate-limiting and abuse monitoring, but the paper questions whether these will scale indefinitely.
\end{itemize}
\end{block}

\begin{columns}[T,totalwidth=\textwidth]
\column{0.30\textwidth}
\begin{block}{Threat actors}
Low/moderate skill actors and advanced persistent threats.
\end{block}

\column{0.30\textwidth}
\begin{block}{Goals}
Profit, chaos, espionage, and specific operational objectives.
\end{block}

\column{0.38\textwidth}
\begin{block}{Misuse applications}
\begin{itemize}
  \item Codex does not yet generate strong standalone malware, but can generate malicious subcomponents.
  \item It can suggest insecure dependencies, insecure function calls, or secrets.
  \item There was limited evidence of major phishing advantages over ordinary language models.
\end{itemize}
\end{block}
\end{columns}

\begin{block}{Data risks}
The slide also notes training-data memorization and data poisoning as concerns when learning from public code corpora.
\end{block}

\refline{\srcChen\sep \srcCarlini\sep \srcBrown}
\end{frame}

% 48
\begin{frame}{Insecure code generation}
\scriptsize
\begin{columns}[T,totalwidth=\textwidth]
\column{0.46\textwidth}
\begin{block}{Generating insecure code}
\begin{itemize}
  \item Because the training corpus is public, Codex may absorb insecure coding practices.
  \item The paper probes cryptographic code generation and measures whether the output is clearly insecure.
\end{itemize}
\end{block}

\column{0.50\textwidth}
\begin{block}{Results placeholder}
\graphplaceholder[0.95\linewidth]{3.0cm}{Insert plot: fraction clearly insecure vs. model size\\for AES and RSA contexts}
\end{block}
\end{columns}

\begin{alertblock}{Criteria used on the slide}
\begin{itemize}
  \item AES contexts are clearly insecure if they use \textbf{ECB mode}.
  \item RSA keys are clearly insecure if they are shorter than \textbf{2048 bits}.
  \item The slide notes that this is probably an underestimate because standards evolve over time.
\end{itemize}
\end{alertblock}

\refline{\srcChen}
\end{frame}

% 49
\begin{frame}{Environmental Impact and Legal Implications (2026 lens)}
\scriptsize
\begin{columns}[T,totalwidth=\textwidth]
\column{0.49\textwidth}
\begin{block}{Environmental impact}
\begin{itemize}
  \item The high-level concern is unchanged: both training and large-scale inference consume meaningful energy.
  \item As coding assistants become always-on and agentic, inference, repeated tool calls, and long-context workflows can dominate operational cost.
  \item The practical response in 2026 is efficiency: better routing, caching, smaller specialist models, and cleaner datacenter power mixes.
\end{itemize}
\end{block}

\column{0.49\textwidth}
\begin{block}{Legal implications}
\begin{itemize}
  \item The 2021 ``fair use + rare copying'' story is no longer sufficient.
  \item In the U.S., the Copyright Office now distinguishes between training-data questions and output-authorship questions.
  \item Output protection depends on meaningful human authorship/control; mere prompting is usually not enough on its own.
  \item In the EU, GPAI obligations make provenance, transparency, and training-content disclosure more operationally important.
\end{itemize}
\end{block}
\end{columns}

{\large\textcolor{black}{\textbf{2026 takeaway}}\par}
\vspace{0.15em}
For deployed coding systems, IP risk management now means provenance, license review, and human accountability - not just hoping the model ``rarely copies''.

\refline{\srcSchwartz\sep \srcIEA\sep \srcUSCOPTwo\sep \srcUSCOPThree\sep \srcEUGPAI}
\end{frame}

% 50
\begin{frame}{Risk Mitigation (2026 lens)}
\scriptsize
\begin{block}{Operational mitigation stack}
\begin{itemize}
  \item Keep agents least-privileged; isolate shell, file-system, network, and secrets access.
  \item Require explicit human approval for high-impact actions: writes, commands, dependency changes, external calls, and pull-request submission.
  \item Log provenance: model version, prompts, tool calls, diffs, test results, and dependency or SBOM changes.
  \item Run automated checks on AI-produced diffs: tests, linting, static analysis, dependency scanning, and secret scanning.
  \item Evaluate on repo-level tasks and red-team for prompt injection, malicious context, and unsafe tool use.
\end{itemize}
\end{block}

{\large\textcolor{black}{\textbf{Principle}}\par}
\vspace{0.25em}
The safe default is to treat AI-generated code and AI-generated actions as \textbf{untrusted until verified}.

\refline{\srcNIST\sep \srcPromptInj\sep \srcOWASP}
\end{frame}

% 51
\outlineframe{9}

% 52
\begin{frame}{Related Work}
\tiny
\begin{columns}[T,totalwidth=\textwidth]
\column{0.49\textwidth}
\begin{block}{Zaremba et al. (2014) -- Learning to Execute}
LSTMs are trained on programs and can learn algorithmic tasks such as adding two 9-digit numbers with about 99\% accuracy. Curriculum learning plays a key role.
\end{block}

\begin{block}{Kulal et al. (2019) -- SPoC}
The SPoC dataset contains about 18K C++ programs and studies pseudocode-to-code synthesis using a seq2seq model with best-first search.
\end{block}

\column{0.49\textwidth}
\begin{block}{Hindle et al. (2012) -- Naturalness of software}
This work studies code with n-gram language models and shows that source code is statistically natural in ways that can be exploited for modeling and tooling.
\end{block}

\begin{block}{Hendrycks et al. (2021) -- APPS}
APPS introduces a benchmark with 10K programming problems across three difficulty levels and includes test cases for evaluating generated solutions.
\end{block}
\end{columns}

\refline{\srcZaremba\sep \srcHindle\sep \srcKulal\sep \srcHendrycks}
\end{frame}

% 53
\begin{frame}{Summary}
\small
\begin{block}{Summary}
\begin{itemize}
  \item The paper studies whether large language models can be trained to generate code from docstrings.
  \item After GitHub fine-tuning, Codex performs well on human-written problems, especially easy interview-style tasks.
  \item Performance improves further when training is brought closer to the evaluation distribution and when multiple samples are considered.
  \item Codex-D is introduced for code-to-docstring generation; it is weaker than code generation but still comparable.
  \item The paper also emphasizes limitations, broader impacts, and safety risks such as over-reliance, bias, misalignment, insecurity, and economic change.
\end{itemize}
\end{block}

\refline{\srcChen}
\end{frame}

% 54
\begin{frame}[plain]
\centering
{\Huge\bfseries The End}
\end{frame}

\end{document}
